% Section 11: Staggered fermions (SF)
% R. Gupta, "Introduction to Lattice QCD", arXiv:hep-lat/9807028
% Transcribed from pages 53-57.
\providecommand{\slashed}[1]{{\mathchoice
  {\ooalign{$\displaystyle#1$\cr\hidewidth$\displaystyle/$\hidewidth}}
  {\ooalign{$\textstyle#1$\cr\hidewidth$\textstyle/$\hidewidth}}
  {\ooalign{$\scriptstyle#1$\cr\hidewidth$\scriptstyle/$\hidewidth}}
  {\ooalign{$\scriptscriptstyle#1$\cr\hidewidth$\scriptscriptstyle/$\hidewidth}}}}

\chapter{Staggered fermions (SF)}

The 16-fold doubling problem of the naive fermion action given in Eq.~7.6
can be reduced to 4 by the trick of spin-diagonalization [21]. Staggered
fermions $\chi$ are defined by the transformation
\begin{equation}
\psi(x) = \Gamma_x \chi(x), \qquad \bar{\psi}(x) = \bar{\chi}(x)\Gamma_x^\dagger,
\qquad \Gamma_x = \gamma_1^{x_1}\gamma_2^{x_2}\gamma_3^{x_3}\gamma_4^{x_4} .
\label{eq:sf-spin-diag}
\end{equation}

In terms of $\chi$ the action can be written as [21]
\begin{equation}
\begin{aligned}
S_S &= m_q\sum_x \bar{\chi}(x)\chi(x)
+ \frac{1}{2}\sum_{x,\mu}\bar{\chi}_x\eta_{x,\mu}
\left[U_{\mu,x}\chi_{x+\hat{\mu}} - U^\dagger_{\mu,x-\hat{\mu}}\chi_{x-\hat{\mu}}\right] \\
&\equiv \sum_{x,y} \bar{\chi}(x) M^S_{xy}\chi(y)
\end{aligned}
\label{eq:sf-action}
\end{equation}
with the matrix $M^S$ given by
\begin{equation}
M^S[U]_{x,y} = m_q\delta_{xy}
+ \frac{1}{2}\sum_\mu\eta_{x,\mu}
\left[U_{x,\mu}\delta_{x,y-\hat{\mu}} - U^\dagger_{x-\hat{\mu},\mu}\delta_{x,y+\hat{\mu}}\right] .
\label{eq:sf-matrix}
\end{equation}

The $\gamma$ matrices are replaced by the phases $\eta_{x,\mu}$. It is convenient to define
the following phase factors
\begin{equation}
\begin{aligned}
\eta_{x,\mu} &= (-1)^{\sum_{\nu<\mu}x_\nu} \\
\zeta_{x,\mu} &= (-1)^{\sum_{\nu>\mu}x_\nu} \\
\epsilon_x &= (-1)^{x_1+x_2+x_3+x_4} \\
S_R(x) &= \frac{1}{2}\left(1 + \zeta_\mu\zeta_\nu - \eta_\mu\eta_\nu
+ \zeta_\mu\zeta_\nu\eta_\mu\eta_\nu\right) .
\end{aligned}
\label{eq:sf-phases}
\end{equation}

From Eqs.~11.2, 11.3 it is easy to see that the different spin components of
$\chi$ are decoupled as the phase factor $\eta_{x,\mu}$ depends only on the site index and
direction and do not have a spinor index. One can therefore drop the spin
index on $\chi$ leaving only color degrees of freedom at each site. This reduces
the original $2^d$-fold degeneracy of naive fermions by a factor of four. The
mass term in Eq.~11.3 is Hermitian, while the $\slashed{D}$ term is anti-Hermitian,
which as I discussed before is required to realize chiral symmetry.

The SF action has translation invariance under shifts by $2a$ due to the
phase factors $\eta_{x,\mu}$. Thus, in the continuum limit, a $2^4$ hypercube is mapped
to a single point, and the 16 degrees of freedom reduce to 4 copies of Dirac
fermions. That is, for every physical flavor, the staggered discretization
produces a 4-fold degeneracy denoted as the staggered flavor. At finite $a$
the gauge interactions break this flavor symmetry and the 16 degrees of
freedom in the hypercube are a mixture of spin and flavor. This is one
of the major drawbacks of staggered fermions. Construction of operators
and their interpretation in terms of spin and flavor is non-trivial as I show
below for the case of bilinears. I begin the discussion of staggered fermions
by listing its symmetries. I will then outline the construction of bilinear
operators and end with a discussion of Ward identities.

\textbf{Symmetries of SF:} The analogue of $\gamma_5$ invariance for SF is
\begin{equation}
\epsilon_x M^{S\dagger}(x,y)\epsilon_y = M^S(x,y),
\qquad \epsilon_x S^{F\dagger}(x,y)\epsilon_y = S^F(y,x) .
\label{eq:sf-gamma5-analog}
\end{equation}
and the symmetries of this action are [63]
\begin{align*}
\text{Translations } S_\mu &: \chi(x) \to \zeta_\mu(x)\chi(x+\hat{\mu}) \\
\text{Rotations } R_{\mu\nu} &: \chi(x) \to S_R(R^{-1}x)\chi(R^{-1}x) \\
\text{Inversion } I &: \chi(x) \to \eta_4(x)\chi(I^{-1}x) \\
\text{Charge Conjugation } C &: \chi(x) \to \epsilon(x)\bar{\chi}(x),
\quad \bar{\chi}(x) \to -\epsilon(x)\chi(x) \\
\text{Vector(Baryon\#) } U(1)_V &: \chi(x) \to e^{i\theta}\chi(x),
\quad \bar{\chi}(x) \to e^{-i\theta}\bar{\chi}(x) \\
\text{Axial } U(1)_A &: \chi(x) \to e^{i\theta\epsilon(x)}\chi(x),
\quad \bar{\chi}(x) \to e^{i\theta\epsilon(x)}\bar{\chi}(x)
\end{align*}
where the phases $\eta$, $\zeta$, $\epsilon$ and $S_R$ are defined in Eq.~11.4. It is the last
invariance, the $U(1)_A$, or also called the $U(1)_\epsilon$ invariance, that becomes a
flavor non-singlet axial symmetry in the continuum limit, and is the main
advantage of staggered fermions.

\textbf{Operators with Staggered fermions:} There are two equivalent ways
of understanding flavor identification of staggered fermions: (i) momentum
space [64], and (ii) position space [65,66]. Here I review the position space
approach.

For the purposes of constructing hadronic operators it is convenient to
represent the 16 degrees of freedom in a hypercube by a hypercube field $Q$,
which in the continuum limit represents the four spin components of four
degenerate flavors [67]. This is done by dividing the lattice into hypercubes
identified by a 4-vector $y$, i.e.\ $y$ is defined on a lattice of spacing $2a$, and
the set of 16 vectors $\{A\}$ to points within the hypercube. These are called
hypercube vectors, and addition of two such vectors is defined modulo 2,
$C_\mu = (A_\mu + B_\mu)\bmod 2$. Then, in the free field limit (setting gauge links
$U_{x,\mu} = 1$ with the understanding that for the interacting theory non-local
operators are to be connected by a path ordered product of links that
makes them gauge invariant, or these operators are to be evaluated in a
fixed gauge)
\begin{equation}
Q(y) = \frac{1}{N_f}\sum_{\{A\}}\Gamma_A\,\chi(y+A)
\label{eq:sf-hypercube-field}
\end{equation}
where $N_f = 4$ and
\begin{equation}
\Gamma_A = \gamma_1^{A_1}\gamma_2^{A_2}\gamma_3^{A_3}\gamma_4^{A_4} .
\label{eq:sf-gamma-A}
\end{equation}

A general bilinear operators can be written as
\begin{equation}
\begin{aligned}
O_{SF} &= \text{Tr}\left[\bar{Q}\gamma_S Q\gamma_F^\dagger\right] \\
&= \sum_{A,B}\bar{\chi}(y+A)\,\text{Tr}\left[\Gamma_A^\dagger\gamma_S\Gamma_B\xi_F\right]\chi(y+B) \\
&= \sum_{A,B}\bar{\chi}(y+A)(\gamma_S\otimes\xi_F)_{AB}\,\chi(y+B) \\
&= \bar{Q}(\gamma_S\otimes\xi_F)Q .
\end{aligned}
\label{eq:sf-bilinear}
\end{equation}
where the matrices $\gamma_S$ determines the spin and $\xi_F$ the flavor of the bilinear.
Each of these can be one of the standard sixteen elements of the Clifford
algebra. In the last expression, the field $Q$ is expressed as a 16 component
vector. It is easy to see from Eq.~11.8 that local operators are given by
$\gamma_S = \xi_F$. The identification of meson operators is as follows [65,68,69].
\begin{equation}
V_\mu = \bar{Q}(\gamma_\mu\otimes\xi_F)Q, \qquad
A_\mu = \bar{Q}(\gamma_\mu\gamma_5\otimes\xi_F)Q, \qquad
P = \bar{Q}(\gamma_5\otimes\xi_F)Q, \qquad
S = \bar{Q}(1\otimes\xi_F)Q .
\label{eq:sf-mesons}
\end{equation}

In the continuum the mesons lie in the flavor SU(4) representations 1 and
15, while on the lattice these representations break into many smaller ones.
For example, the continuum 15-plet of pions breaks into seven lattice representations (four 3-d and three 1-d).
\begin{equation}
\begin{aligned}
15 &\to (\gamma_5\otimes\gamma_i) \oplus (\gamma_5\otimes\gamma_i\gamma_4)
\oplus (\gamma_5\otimes\gamma_i\gamma_5) \oplus (\gamma_5\otimes\gamma_i\gamma_4\gamma_5) \\
&\quad\oplus (\gamma_5\otimes\gamma_4) \oplus (\gamma_5\otimes\gamma_5) \oplus (\gamma_5\otimes\gamma_4\gamma_5), \\
1 &\to (\gamma_5\otimes 1) .
\end{aligned}
\label{eq:sf-pion-reps}
\end{equation}
In the first line, $i = 1, 2, 3$, so the representations are 3-dimensional. The
other representations are 1-dimensional. Similarly, the continuum rho 15-plet (times three spin components) breaks into eleven representations (four
6-d and seven 3-d).
\begin{equation}
\begin{aligned}
15 &\to (\gamma_i\otimes\gamma_j) \oplus (\gamma_i\otimes\gamma_j\gamma_4)
\oplus (\gamma_i\otimes\gamma_j\gamma_5) \oplus (\gamma_i\otimes\gamma_j\gamma_4\gamma_5) \\
&\quad\oplus (\gamma_i\otimes\gamma_i) \oplus (\gamma_i\otimes\gamma_i\gamma_4)
\oplus (\gamma_i\otimes\gamma_i\gamma_5) \oplus (\gamma_i\otimes\gamma_i\gamma_4\gamma_5) \\
&\quad\oplus (\gamma_i\otimes\gamma_4) \oplus (\gamma_i\otimes\gamma_5) \oplus (\gamma_i\otimes\gamma_4\gamma_5), \\
1 &\to (\gamma_i\otimes 1) .
\end{aligned}
\label{eq:sf-rho-reps}
\end{equation}

Of the sixteen pions, the Goldstone pion has flavor $\xi_5$. The remaining 15
pions are given by the other 15 members of the Clifford algebra and are
heavier due to the staggered flavor symmetry breaking, and their masses
do not vanish at $m_q = 0$. In fact the mass for each representation can be
different. The degree to which the staggered flavor symmetry is broken can
be quantified, for example, by comparing the mass of the Goldstone pion
versus those in the other representations.

The above examples make it clear that the mixing of spin and flavor
complicates the analysis and is a major disadvantage of staggered fermions.
I will not go into details of calculations using staggered fermions. A useful
list of references is as follows. The construction of baryon operators is
given in [70]. For an illustrative calculations of the spectrum see [69]. The
transcription of 4-fermion operators arising in effective weak Hamiltonian
and their matrix elements within hadronic states are discussed in [63,71].

\textbf{Chiral Symmetry and Ward Identities:} The staggered fermion action
can be written in terms of the 16 component fields $Q$, defined in Eq.~11.6,
as [4]
\begin{equation}
S_S = m_q\sum_y \bar{Q}(y)(1\otimes 1)Q(y)
+ \sum_{y,\mu}\bar{Q}(y)\left[(\gamma_\mu\otimes 1)\partial_\mu
+ (\gamma_5\otimes\xi_\mu\xi_5)\right]_\mu Q(y) .
\label{eq:sf-action-hypercube}
\end{equation}
This action, for $m_q = 0$, is invariant under the following abelian transformation
\begin{equation}
Q(y) \to e^{i\theta(\gamma_5\otimes\xi_5)}Q(y), \qquad
\bar{Q}(y) \to \bar{Q}(y)e^{i\theta(\gamma_5\otimes\xi_5)} .
\label{eq:sf-abelian}
\end{equation}
Thus one preserves a non-trivial piece of the original $U(4)\times U(4)$ symmetry
(called $U(1)_\epsilon$ or $U(1)_A$) whose generator is $\gamma_5\otimes\xi_5$ and the associated
Goldstone boson operator is $\bar{Q}(\gamma_5\otimes\xi_5)Q$. This axial $U(1)_A$ symmetry,
along with the $U(1)_V$, allows, for appropriate transcription of operators,
Ward Identities for lattice amplitudes similar to those in the continuum.

The best studied example of this is the lattice calculation of the $K^0\bar{K}^0$
mixing parameter $B_K$ [72].

The basic tools for constructing the Ward Identities are exemplified by
the following two relations involving quark propagators and the zero 4-momentum insertion of the pseudo-Goldstone operator $\bar{\chi}(x)\chi(x)(-1)^x$ and
the scalar density $\bar{\chi}(x)\chi(x)$ [63]
\begin{equation}
(-1)^{x_f}G_2(x_f,x_i) + G_1(x_f,x_i)(-1)^{x_i}
= (m_1+m_2)\sum_x G_1(x_f,x)(-1)^x G_2(x,x_i)
\label{eq:sf-wi-1}
\end{equation}
\begin{equation}
G_2(x_f,x_i) - G_1(x_f,x_i)
= (m_1-m_2)\sum_x G_1(x_f,x)G_2(x,x_i)
\label{eq:sf-wi-2}
\end{equation}
where $G_i(x_f,x_i)$ is the propagator from $x_i$ to $x_f$ for a quark of mass of
mass $m_i$. The derivation of these identities, using the hopping parameter
expansion, is given in Ref.~[63]. Note that these relations are valid for all
$a$ and at finite $m$. An example of the usefulness of these relations is in
extracting the chiral condensate as shown below. Eq.~11.14 gives
\begin{equation}
\langle\bar{\chi}\chi\rangle = G(x_i,x_i)
= m\sum_x (-1)^{-x_i}G(x_i,x)(-1)^x G(x,x_i)
= m\sum_x |G(x_i,x)|^2 ,
\label{eq:sf-condensate}
\end{equation}
where the last term on the right is the zero 4-momentum sum of the pion
correlator. Similarly Eq.~11.15 gives
\begin{equation}
\frac{\partial\langle\bar{\chi}\chi\rangle}{\partial m}
= \sum_x G(x_i,x)G(x,x_i)
= \sum_x |G(x_i,x)|^2(-1)^{x-x_i} .
\label{eq:sf-condensate-deriv}
\end{equation}
Combining these two equations one can extract the value of the condensate
in the chiral limit
\begin{equation}
\langle\bar{\chi}\chi\rangle(m=0)
= \left(1 - m\frac{\partial}{\partial m}\right)\langle\bar{\chi}\chi\rangle
= \sum_{\mathrm{even}} |G_1(x_i,x)|^2
\label{eq:sf-condensate-chiral}
\end{equation}
by summing the Goldstone pion's 2-point function over sites with even
$(x-x_i)$.

\section{Summary}

Even though most of the simulations done until recently have used either
the staggered or Wilson formulation of fermions, neither is satisfactory.
